% Draft prose for the remaining gaps in "Higher-dimensional realistic
% data" and the "Conclusion" section. Ready to paste into main.tex; not
% wired in itself.
%
% Revision note: this version corrects a real error in the previous
% draft. The wind-speed paragraph there implied the MDN rediscovers the
% curtailed-vs-on-curve bimodality the EDA hypothesized (Figure
% fig:wind-conditional-histograms). Checked directly against the trained
% models: it does not. The "genuine" well-weighted splits found at 9 m/s
% (K=4) and 13 m/s (K=3) are BOTH components sitting on/near the main
% trend line, not one curtailed and one on-curve. The actual curtailed
% mode is real in the data (verified independent of any model) but is
% represented with far less weight than its true frequency, at every K
% and speed checked except 13 m/s, where curtailment is rare enough that
% the mismatch mostly disappears. This is now stated explicitly rather
% than glossed over.

% ============================================================
% 1. Curve overlay + weight sweep INTERPRETATION
%    (goes after the existing weight-definition paragraph, line 292,
%    and before \reportfigurepairtall{curve_overlay_speed...}, line 296)
% ============================================================

Wind speed provides the clearest evidence that the MDN discovers real secondary structure, but it is not the structure our exploratory analysis hypothesized.
Figure~\ref{fig:wind-conditional-histograms} motivated wind speed as multimodal because a fixed-speed slice of the real data shows two clusters: a curtailed, near-zero-power cluster and an on-curve cluster.
Checking this directly against the trained models, that curtailed mode is represented with far less weight than its true frequency at every $K$ and speed we checked, with one exception.
At 9 m/s, where 6.7\% of real data points fall below 200 kW, the only component near 0 kW across $K=2$ through $K=5$ is $K=4$'s orange component (mean $-16.5$ kW), carrying just 0.13\% of the mixture weight.
At 5 m/s, where curtailment is far more common (19.6\% of points below 200 kW), the near-zero component is $K=3$'s green component (mean 10 kW, 2.3\% weight), $K=4$'s orange component (mean $-10$ kW, 2.4\% weight), and $K=5$'s orange component (mean $-0.1$ kW, 1.3\% weight) -- roughly an order of magnitude below the true rate in every case, though genuinely present near zero, unlike at 9 m/s.
Only at 13 m/s, where curtailment has dropped to 1.5\% of points, does the near-absence of a near-zero component look reasonable rather than wrong.

What the MDN instead represents with substantial, well-balanced weight is a second \emph{on-curve} mode, not a curtailed one.
At $K=4$ and 9 m/s, 59\% of the weight sits on the red component (mean near 1818 kW) and 41\% on the green component (mean near 1917 kW) -- both close to the main trend line, not one curtailed and one on-curve.
Because these two means are only 99 kW apart relative to an axis spanning several thousand kW, the two lines sit almost on top of each other in Figure~\ref{fig:wind-curve-overlay-speed-direction} at this point -- visually subtle, but confirmed as substantial by the weight panel alongside it.
That weight panel also shows this is not an isolated point: across most of the swept range a different color is dominant in turn (orange below roughly 3.5 m/s, red through the early-to-mid ramp, green through the rest of the ramp, blue above roughly 13 m/s).
$K=4$ carries the highest average secondary weight of any $K$ across the full speed sweep (0.138), and $K=3$ independently recovers a similarly on-trend split near rated wind speed: at 13 m/s, $K=3$ splits 80/20 between the blue component (3526 kW) and the orange component (3396 kW).

We read this as a genuine, if unexpected, finding rather than a failure to reproduce the exploratory analysis: whatever secondary structure the MDN finds useful enough to spend real weight on is not the curtailment split we hypothesized going in, but a smaller, on-trend split whose physical cause we have not identified.
The curtailment mode is real -- we verified it directly in the raw data, independent of any trained model -- but the MDN does not represent it with weight proportional to its actual frequency outside the regime near rated power, where curtailment itself becomes rare.

Wind direction and day-of-year present a substantially weaker case (Figure~\ref{fig:wind-curve-overlay-doy}).
In the wind-direction sweep (left), how the diverging component behaves depends on $K$: at $K=4$ and $K=5$, the diverging component (blue and red respectively) dips sharply within a narrow slice of the domain, roughly $75$-$100^\circ$, and recovers toward the main curve elsewhere, while at $K=2$ and $K=3$ the diverging component (blue and green respectively) instead stays low across nearly the entire domain rather than in one narrow slice.
Neither behavior tracks a real competing mode across most of the feature's range the way wind speed's components do.
The day-of-year sweep (right) shows the dominant component (orange at $K=2,3$) elevated near both the start and end of the calendar year (roughly days 0-50 and 300-365), with a lower, flatter plateau in between, rather than splitting into two persistently co-present modes; as noted earlier, this pattern is confounded with a single year's weather and cannot be taken as evidence of a genuine repeating seasonal effect.
Both features' apparent secondary weight is also far less stable across $K$ than wind speed's: wind direction's average secondary weight jumps from 0.04-0.07 at $K=2,3$ to 0.27 at $K=4$ before dropping back to 0.17 at $K=5$, and day-of-year shows the same $K=4$-concentrated pattern (0.06-0.08 at $K=2,3$ versus 0.21 at $K=4$).
A pattern that only appears at one specific $K$ is weaker evidence than one that persists across several, since it is consistent with a particular optimization run happening to find a use for extra capacity rather than the model consistently rediscovering real structure regardless of how many components it is given -- and this comparison rests on a single training seed, so we cannot rule out that the $K=4$ concentration is itself an artifact of this particular run rather than a reproducible property of the data.
Taken together, direction and day-of-year contribute real but inconsistent signal that we would not treat as evidence of genuine multimodal structure with the same confidence as wind speed's on-trend split -- and, as with wind speed, we have not confirmed that either feature's signal corresponds to the curtailment-like pattern originally hypothesized in the EDA rather than some other structure.

% ============================================================
% 2. Density grid - short explanation
%    (goes after \reportfigure{density_grid_wind_turbine.png}, line 309)
% ============================================================

Figure~\ref{fig:density-grid-wind} shows the full predicted density $p(\text{ActivePower} \mid \mathbf{x})$ at three representative wind-speed values (5, 9, and 13 m/s, spanning the ramp region and rated output) across all five values of $K$, complementing the curve-overlay sweep by showing the actual shape of the predicted distribution at fixed points rather than only its swept mean and weight.
This view makes the mismatch described above directly visible: none of the density panels show meaningful mass near 0 kW at 5 or 9 m/s, despite 19.6\% and 6.7\% of the real data sitting below 200 kW at those speeds respectively -- the predicted densities are concentrated almost entirely around the on-curve value, with at most a thin, low-height tail toward zero.
The splits that \emph{are} visible as genuinely separated probability mass are the on-trend splits discussed above, not a curtailed-vs-on-curve split: at $K=4$ and 9 m/s, two distinct peaks near 1818 kW and 1917 kW, and at $K=3$ and 13 m/s, near 3396 kW and 3526 kW.
In both cases the peaks visibly overlap given each component's spread, meaning even this genuine structure is not always sharply distinguishable point-to-point -- a point relevant to the model's predictive sharpness and calibration, separate from the question of whether it represents curtailment at all.

The grid also reinforces the point made above by direct visual contrast.
$K=1$ always shows a single gaussian by construction, but $K=2$ collapses to an effectively single dominant mode at all three wind speeds (98\%+ weight on one component at 5 and 9 m/s, and 99.99\% at 13 m/s) despite having a second component available, and even $K=3$ through $K=5$ spend most of their extra components on the on-trend split rather than on the curtailment mode the raw data actually contains.
Simply allowing more components does not guarantee the model uses them for the structure we expected -- it may use them for something else entirely, as is the case here.

% ============================================================
% 3. Conclusion
%    (replaces the answerbox placeholder, lines 338-340)
%    Simplified per request: focus on regression-vs-MDN performance and
%    how the MDN did on the synthetic task vs. the wind-turbine task.
%    Numbers pulled from the current main.tex (lines 177, 255, 261).
% ============================================================

On the synthetic task, the MDN beats the deterministic regression model outright: every MDN, including $K=1$, achieves better test RMSE than RegressionNet, and $K=2$ gives the best result on both metrics (test RMSE 0.9036, test NLL $-0.2221$), matching the true two-mode generative process.
Here, representing the two modes explicitly helped point accuracy as well as distributional fit -- there was no trade-off to make.

On the wind-turbine task, the two model families split the comparison instead of one simply winning.
RegressionNet achieves the best test RMSE overall (195.93 kW), beating every MDN including the best of them ($K=5$, 223.37 kW), while the MDNs still achieve better test NLL than the single-Gaussian model as $K$ increases.
This is not a contradiction: RegressionNet is trained to directly minimize MSE and so converges toward the conditional mean, the MSE-optimal point predictor, while an MDN's point estimate is only a byproduct of maximizing likelihood.
Where the MDN represents two components as genuinely separated, the mixture mean falls between them in a region of low density, giving a worse point estimate even though the underlying distribution is better represented.

So the MDN's advantage over regression is consistently a distributional one (NLL), not a point-accuracy one (RMSE) -- but whether that costs anything on RMSE depends on the task.
On the synthetic data, the two modes were separated enough that fitting them explicitly improved point accuracy too.
On the wind-turbine data, real separation between components exists but comes at a measurable RMSE cost, so RegressionNet remains the better choice if a single point prediction is all that is needed, while the MDN is the better choice if the goal is representing the actual spread of likely outcomes.
